\documentclass[10pt,a4paper,oneside]{book}

% -------------------------------------------------
% LINGUA, FONT E MATEMATICA
% -------------------------------------------------
\usepackage[italian]{babel}
\usepackage[utf8]{inputenc}
\usepackage[T1]{fontenc}
\usepackage{lmodern}
\usepackage{microtype}
\usepackage{amsmath,amssymb,mathtools}

% -------------------------------------------------
% IMPAGINAZIONE
% -------------------------------------------------
\usepackage[
    a4paper,
    top=1.55cm,
    bottom=1.60cm,
    left=1.55cm,
    right=1.55cm,
    headsep=0.35cm,
    footskip=0.80cm
]{geometry}

\usepackage{setspace}
\setstretch{1.05}
\usepackage{parskip}
\setlength{\parskip}{5pt}
\setlength{\parindent}{0pt}
\usepackage{enumitem}
\setlist[itemize]{topsep=2pt,itemsep=2pt}

% -------------------------------------------------
% TITOLI
% -------------------------------------------------
\usepackage{titlesec}
\titleformat{\chapter}{\normalfont\LARGE\bfseries}{\thechapter.}{0.6em}{}
\titleformat{\section}{\normalfont\Large\bfseries}{\thesection}{0.6em}{}
\titleformat{\subsection}{\normalfont\large\bfseries}{\thesubsection}{0.6em}{}
\titlespacing*{\chapter}{0pt}{-4pt}{12pt}
\titlespacing*{\section}{0pt}{12pt}{5pt}
\titlespacing*{\subsection}{0pt}{8pt}{3pt}

% -------------------------------------------------
% BOX COLORATI
% -------------------------------------------------
\usepackage{xcolor}
\usepackage[most]{tcolorbox}

\tcbset{
    enhanced,
    breakable,
    boxrule=0.4pt,
    arc=1mm,
    left=6pt,
    right=6pt,
    top=5pt,
    bottom=5pt,
    before skip=7pt,
    after skip=7pt
}

\newtcolorbox{definitionbox}[1]{
    title=\textbf{Definizione -- #1},
    colback=blue!4,
    colframe=blue!55!black,
    colbacktitle=blue!55!black,
    coltitle=white
}

\newtcolorbox{importantbox}[1]{
    title=\textbf{Idea chiave -- #1},
    colback=orange!7,
    colframe=orange!70!black,
    colbacktitle=orange!70!black,
    coltitle=white
}

\newtcolorbox{examplebox}[1]{
    title=\textbf{Esempio -- #1},
    colback=green!5,
    colframe=green!45!black,
    colbacktitle=green!45!black,
    coltitle=white
}

\newtcolorbox{warningbox}[1]{
    title=\textbf{Attenzione -- #1},
    colback=red!5,
    colframe=red!65!black,
    colbacktitle=red!65!black,
    coltitle=white
}

% -------------------------------------------------
% INTESTAZIONE E LINK
% -------------------------------------------------
\usepackage{fancyhdr}
\pagestyle{fancy}
\fancyhf{}
\fancyhead[L]{XAI Counterfactuals}
\fancyhead[R]{\nouppercase{\leftmark}}
\fancyfoot[C]{\thepage}
\setlength{\headheight}{13pt}
\usepackage[hidelinks]{hyperref}

% -------------------------------------------------
% FRONTESPIZIO
% -------------------------------------------------
\title{\textbf{XAI Counterfactuals}\\[0.4em]
       \large Quaderno di progetto e indice dei file}
\author{}
\date{\today}

\begin{document}

\maketitle
\tableofcontents

% =================================================
\chapter{Guida al progetto}
% =================================================

\section{Obiettivo e impostazione}

Il progetto studia spiegazioni controfattuali per un classificatore di immagini
\textit{real/fake}. La prima architettura prevista è una ResNet18 pre-addestrata
in PyTorch; in seguito verranno considerati i quattro metodi controfattuali
disponibili nella libreria Alibi. Un'eventuale estensione a Vision Transformer
(ViT) dipenderà dai risultati dei primi esperimenti.

\begin{importantbox}{Finalità di questo documento}
Questo file è un riferimento interno per ricordare \textbf{a cosa serve ogni
componente della repository} e raccogliere progressivamente note su
implementazione, esperimenti e decisioni. Non sostituisce il
\texttt{WORKLOG.md}, che registra l'avanzamento cronologico del lavoro.
\end{importantbox}

\section{Struttura della repository}

\begin{verbatim}
xai-counterfactuals/
|-- main.py
|-- requirements.txt
|-- README.md
|-- WORKLOG.md
|-- .gitignore
|-- src/
|   |-- __init__.py
|   |-- config.py
|   |-- dataset.py
|   |-- models.py
|   |-- train.py
|   |-- evaluate.py
|   `-- explainers/
|       |-- __init__.py
|       |-- cf.py
|       |-- cem.py
|       |-- cfproto.py
|       `-- cfrl.py
|-- experiments/
|   `-- 01_exploration.ipynb
|-- data/
|   |-- real/
|   `-- fake/
|-- checkpoints/
|-- results/
|-- slides/
|   `-- presentation.tex
`-- docs/
    `-- project_reference.tex
\end{verbatim}

\begin{warningbox}{Stato dei file}
La struttura qui rappresentata è quella \textbf{prevista}. Alcuni file possono
essere ancora vuoti o non essere stati salvati nella repository. L'indice
va aggiornato man mano che l'implementazione procede.
\end{warningbox}


\section{File nella directory principale}

\subsection{\texttt{main.py}}

Punto di ingresso previsto per avviare i principali esperimenti.

\subsection{\texttt{requirements.txt}}

Elenco delle dipendenze Python necessarie al progetto.

\subsection{\texttt{README.md}}

Presentazione della repository, obiettivi e istruzioni d'uso.

\subsection{\texttt{WORKLOG.md}}

Diario settimanale delle attività, dei problemi e delle decisioni.

\subsection{\texttt{.gitignore}}

Esclude dal controllo di versione dati, checkpoint e file temporanei.


\section{Moduli Python: \texttt{src/}}

\subsection{\texttt{config.py}}

Percorsi del progetto, parametri di training e scelta del dispositivo
(CPU, MPS o CUDA).


\subsection{\texttt{dataset.py}}

Il file \texttt{dataset.py} si occupa di caricare le immagini del dataset,
applicare le trasformazioni necessarie per ResNet18 e organizzare le
osservazioni in batch per l'addestramento e la valutazione.

Il codice è organizzato in due funzioni principali:
\texttt{get\_transforms()} e \texttt{get\_dataloaders()}.

\subsubsection{Preprocessing delle immagini}

La funzione \texttt{get\_transforms()} definisce una sequenza di
trasformazioni attraverso \texttt{transforms.Compose()}:

\begin{verbatim}
transform = transforms.Compose([
    transforms.Resize((IMAGE_SIZE, IMAGE_SIZE)),
    transforms.ToTensor(),
    transforms.Normalize(
        mean=[0.485, 0.456, 0.406],
        std=[0.229, 0.224, 0.225]
    )
])
\end{verbatim}

\textbf{transforms.Compose()} permette di concatenare più trasformazioni,
che vengono applicate nell'ordine specificato a ogni immagine.

Le tre operazioni sono:

\begin{enumerate}
    \item \textbf{Resize}: ridimensiona l'immagine a
    $224 \times 224$ pixel, indipendentemente dalla sua dimensione
    originale.

    \item \textbf{ToTensor}: converte l'immagine in un tensore PyTorch.
    Per immagini RGB con valori interi tra 0 e 255, la trasformazione
    converte i pixel in valori floating-point nell'intervallo $[0,1]$.
    Inoltre, cambia l'ordine delle dimensioni da
    $(H,W,C)$ a $(C,H,W)$.

    \item \textbf{Normalize}: standardizza separatamente i tre canali
    RGB utilizzando le medie e le deviazioni standard di ImageNet.
\end{enumerate}

La normalizzazione viene effettuata per ogni pixel del canale $c$:

\[
z_{c,h,w}
=
\frac{x_{c,h,w}-\mu_c}{\sigma_c},
\]

dove $\mu_c$ e $\sigma_c$ sono rispettivamente la media e la deviazione
standard utilizzate per il canale $c$.

\begin{importantbox}{Perché utilizziamo i parametri di ImageNet?}
ResNet18 è pre-addestrata su ImageNet. Utilizzare la stessa
normalizzazione permette di fornire al modello immagini preprocessate
in modo coerente con quelle utilizzate durante il pre-training.

Questa trasformazione non garantisce che il dataset FFHQ abbia media
zero e varianza unitaria: utilizziamo statistiche fissate, calcolate
su ImageNet.
\end{importantbox}

Al termine del preprocessing, una singola immagine è rappresentata
da un tensore:

\[
x_i \in \mathbb{R}^{3 \times 224 \times 224}.
\]

\subsubsection{Caricamento del dataset con ImageFolder}

Il caricamento delle immagini viene effettuato attraverso:

\begin{verbatim}
dataset = datasets.ImageFolder(
    root=DATA_DIR,
    transform=get_transforms()
)
\end{verbatim}

\textbf{ImageFolder} è una classe di Torchvision che costruisce
automaticamente un dataset a partire da una directory organizzata
in sottocartelle, ciascuna corrispondente a una classe (https://docs.pytorch.org/vision/main/generated/torchvision.datasets.ImageFolder.html). 


Nel nostro caso:

\begin{verbatim}
data/
    fake/
    real/
\end{verbatim}

Ogni immagine viene associata a un'etichetta numerica secondo
l'ordine alfabetico delle sottocartelle:

\[
y_i =
\begin{cases}
0 & \text{se l'immagine è fake},\\
1 & \text{se l'immagine è real}.
\end{cases}
\]

L'oggetto \texttt{dataset} permette di accedere alle osservazioni
tramite un indice. Ad esempio, \texttt{dataset[i]} restituisce
una coppia:

\[
(x_i,y_i),
\]

dove $x_i$ è il tensore dell'immagine dopo il preprocessing
e $y_i$ è la sua etichetta.

È importante osservare che \texttt{ImageFolder} non carica
necessariamente tutte le immagini in memoria: mantiene i percorsi
dei file e legge le immagini quando vengono richieste.

\subsubsection{Suddivisione in training, validation e test}

Il dataset viene suddiviso utilizzando la funzione
\texttt{random\_split()} di PyTorch:

\begin{verbatim}
generator = torch.Generator().manual_seed(RANDOM_SEED)

train_set, val_set, test_set = random_split(
    dataset,
    [n_train, n_val, n_test],
    generator=generator
)
\end{verbatim}

La funzione \textbf{random\_split()} genera una partizione casuale
degli indici del dataset, suddividendo le osservazioni in tre
sottoinsiemi non sovrapposti.

Le dimensioni sono definite da:

\[
n_{\mathrm{train}} = \lfloor 0.70n \rfloor,
\]

\[
n_{\mathrm{val}} = \lfloor 0.15n \rfloor,
\]

\[
n_{\mathrm{test}}
=
n-n_{\mathrm{train}}-n_{\mathrm{val}}.
\]

Il training set viene utilizzato per stimare i parametri del modello,
il validation set per monitorare e selezionare il modello durante
lo sviluppo, mentre il test set viene riservato alla valutazione finale.

\begin{definitionbox}{torch.Generator e manual\_seed}
\texttt{torch.Generator()} crea un generatore di numeri
pseudo-casuali. Il metodo \texttt{manual\_seed()} ne fissa il seme.

Utilizzando lo stesso seed, \texttt{random\_split()} produce
la stessa suddivisione, a condizione che il dataset e il suo ordine
rimangano invariati.

Questo è importante per la riproducibilità degli esperimenti.
\end{definitionbox}

\subsubsection{DataLoader e costruzione dei batch}

Dopo la suddivisione, ciascun sottoinsieme viene passato a un
\texttt{DataLoader}:

\begin{verbatim}
train_loader = DataLoader(
    train_set,
    batch_size=BATCH_SIZE,
    shuffle=True
)
\end{verbatim}

Il \textbf{DataLoader} permette di iterare sul dataset attraverso
gruppi di osservazioni, chiamati \textit{mini-batch}, invece di
elaborare singolarmente ogni immagine.


Con \texttt{BATCH\_SIZE = 32}, il training set viene suddiviso
in mini-batch contenenti al massimo 32 osservazioni ciascuno.

Indicando con $n_{\mathrm{train}}$ il numero di osservazioni
nel training set, il numero totale di batch per epoca è:

\[
B = \left\lceil \frac{n_{\mathrm{train}}}{32} \right\rceil,
\]

dove $\lceil \cdot \rceil$ rappresenta la funzione ceiling,
che arrotonda all'intero superiore.

Indichiamo con $b \in \{1,\ldots,B\}$ l'indice del batch.
Per un batch completo, i dati e le rispettive etichette
hanno dimensioni:

\[
X^{(b)} \in \mathbb{R}^{32 \times 3 \times 224 \times 224},
\qquad
y^{(b)} \in \{0,1\}^{32}.
\]

In generale, indicando con $m_b$ il numero di osservazioni
del batch $b$, abbiamo:

\[
X^{(b)} \in \mathbb{R}^{m_b \times 3 \times 224 \times 224},
\qquad
y^{(b)} \in \{0,1\}^{m_b},
\]

con $1 \leq m_b \leq 32$. L'ultimo batch può infatti
contenere meno di 32 osservazioni.

\begin{definitionbox}{Epoca (Epoch)}
Un'\textbf{epoca} corrisponde a un passaggio completo
attraverso tutte le osservazioni del training set.

Durante un'epoca, il modello elabora sequenzialmente
i $B$ mini-batch. Per ciascun batch vengono calcolati
le predizioni, la funzione di perdita e, durante il training,
l'aggiornamento dei parametri del modello.

Pertanto, con un aggiornamento per batch, un'epoca
comprende $B$ iterazioni di ottimizzazione.
\end{definitionbox}

\begin{examplebox}{Numero di batch per epoca}
Supponiamo che il training set contenga
$n_{\mathrm{train}}=100$ immagini e che
\texttt{BATCH\_SIZE = 32}.

Il numero di batch sarà:

\[
B = \left\lceil \frac{100}{32} \right\rceil = 4.
\]

Durante un'epoca il modello elaborerà:
\[
m_1=32,\quad m_2=32,\quad m_3=32,\quad m_4=4.
\]

Al termine dei quattro batch, tutte le 100 osservazioni
saranno state utilizzate una volta e l'epoca sarà conclusa.
\end{examplebox}

Il parametro \texttt{shuffle=True} fa sì che, all'inizio
di ogni nuova epoca, venga generato un nuovo ordinamento
casuale delle osservazioni del training set.

Di conseguenza, ogni epoca utilizza le stesse osservazioni,
ma generalmente raggruppate in mini-batch differenti.

Nei DataLoader di validation e test utilizziamo invece
\texttt{shuffle=False}, perché non è necessario modificare
l'ordine delle osservazioni per valutare il modello.

\begin{importantbox}{Dataset e DataLoader}
Il \texttt{Dataset} definisce come accedere a una singola
osservazione $(x_i,y_i)$.

Il \texttt{DataLoader} organizza queste osservazioni in batch,
gestendone l'ordine e il caricamento durante l'esecuzione.

Quindi il Dataset rappresenta la sorgente dei dati, mentre
il DataLoader definisce come i dati vengono forniti al modello.
\end{importantbox}

\begin{warningbox}{Limiti dell'implementazione iniziale}
La funzione \texttt{random\_split()} non garantisce una
suddivisione stratificata rispetto alle classi. Pertanto,
anche partendo da un dataset bilanciato, i singoli sottoinsiemi
potrebbero contenere proporzioni diverse di immagini real e fake.

Inoltre, tutte le immagini utilizzano attualmente lo stesso
preprocessing, senza trasformazioni aggiuntive specifiche per
il training.

Queste scelte sono sufficienti per verificare inizialmente
la pipeline, ma dovranno essere rivalutate prima degli
esperimenti definitivi.
\end{warningbox}

\begin{warningbox}{Il significato di \texttt{shuffle=True}}
Nel \texttt{DataLoader} di PyTorch, il parametro
\texttt{shuffle=True} non indica un campionamento casuale
con reinserimento, ma una \textbf{permutazione casuale
senza reinserimento} degli indici del dataset.

All'inizio di ogni epoca viene generata una permutazione
casuale $\pi$ degli $n_{\mathrm{train}}$ indici. Questa
permutazione viene successivamente suddivisa in $B$
mini-batch consecutivi.

Di conseguenza, i mini-batch costituiscono una partizione
degli indici del training set:

\[
\bigcup_{b=1}^{B}\mathcal{B}_b
=
\{1,\ldots,n_{\mathrm{train}}\},
\]

\[
\mathcal{B}_i \cap \mathcal{B}_j
=
\varnothing
\qquad \forall i\neq j.
\]

Questo garantisce che \textbf{ogni osservazione venga
utilizzata esattamente una volta per epoca}, anche se
l'ordine e la composizione dei mini-batch cambiano
casualmente tra epoche successive.

La proprietà vale con \texttt{drop\_last=False}
(impostazione predefinita). Se invece
\texttt{drop\_last=True}, l'ultimo mini-batch incompleto
viene scartato e alcune osservazioni possono non essere
utilizzate durante quella specifica epoca.
\end{warningbox}



\subsection{\texttt{models.py}}

Il file \texttt{models.py} definisce il classificatore utilizzato
nel progetto. Inizialmente viene impiegata una \textbf{ResNet18},
una rete neurale convoluzionale pre-addestrata su ImageNet,
successivamente adattata al problema di classificazione binaria
tra immagini \textit{real} e \textit{fake}.

Non implementiamo manualmente l'architettura della rete, poiché
Torchvision fornisce già una ResNet18 completa, comprensiva
dei layer convoluzionali, dei blocchi residuali e dei relativi
pesi pre-addestrati. Il codice definisce un'unica classe, \texttt{ResNet18Classifier},
che eredita da \texttt{nn.Module}.

\subsubsection{Modelli pre-addestrati e transfer learning}

La funzione \texttt{resnet18()} di Torchvision permette di
istanziare l'architettura ResNet18 scegliendo se utilizzare
pesi pre-addestrati oppure inizializzazioni casuali.

\begin{verbatim}
model = resnet18(weights=None)
\end{verbatim}

In questo caso viene costruita l'architettura, ma i parametri
non provengono da un precedente addestramento. Nel nostro progetto utilizziamo invece:

\begin{verbatim}
model = resnet18(
    weights=ResNet18_Weights.DEFAULT
)
\end{verbatim}

L'argomento \texttt{weights} specifica quali pesi pre-addestrati
caricare. Con \texttt{ResNet18\_Weights.DEFAULT} vengono
utilizzati i pesi pre-addestrati di riferimento forniti
da Torchvision per ImageNet-1K, un problema di classificazione con 1000 classi.

\begin{definitionbox}{Transfer Learning}
Il \textbf{Transfer Learning} consiste nel riutilizzare
conoscenze apprese da un modello durante un precedente task
per affrontare un nuovo problema.

Indicando con $\theta_{\mathrm{pre}}$ i parametri ottenuti
durante il pre-training, invece di partire da una
inizializzazione casuale $\theta_0$, utilizziamo:

\[
\theta_0 = \theta_{\mathrm{pre}},
\]

per i layer di cui riutilizziamo i pesi.

Nel nostro caso, le rappresentazioni visive apprese
su ImageNet vengono riutilizzate per la classificazione
di volti reali e generati.
\end{definitionbox}

Questo approccio risulta particolarmente utile quando
il dataset disponibile per il nuovo task è relativamente
piccolo. Infatti, i layer della rete hanno già appreso
rappresentazioni utili, come bordi, texture e strutture
visive più complesse.

Il modello necessita comunque di essere adattato al nuovo
problema: i pesi pre-addestrati non permettono direttamente
di distinguere immagini reali da immagini generate.

\subsubsection{La classe ResNet18Classifier}

La classe viene dichiarata come:

\begin{verbatim}
class ResNet18Classifier(nn.Module):
\end{verbatim}

\texttt{nn.Module} è la classe base di PyTorch per la
costruzione di reti neurali. Ereditare da questa classe
permette di utilizzare automaticamente funzionalità
come la gestione dei parametri, il trasferimento del
modello su GPU e le modalità di training e valutazione.

All'interno del costruttore troviamo:

\begin{verbatim}
def __init__(self):
    super().__init__()
\end{verbatim}

La chiamata \texttt{super().\_\_init\_\_()} inizializza
la classe padre \texttt{nn.Module}, rendendo possibile
la registrazione dei layer e dei parametri del modello.

Successivamente carichiamo la ResNet18:

\begin{verbatim}
self.model = resnet18(
    weights=ResNet18_Weights.DEFAULT
)
\end{verbatim}

L'attributo \texttt{self.model} contiene quindi
l'intera rete neurale pre-addestrata.

\subsubsection{Adattamento del classificatore finale}

La ResNet18 originale è stata addestrata per distinguere
1000 classi di ImageNet. Il suo ultimo layer è un
\textit{fully connected layer}, accessibile tramite
l'attributo \texttt{fc}.

Questo layer riceve un vettore di 512 caratteristiche
estratte dalla rete e restituisce 1000 valori:

\[
\mathbb{R}^{512}
\longrightarrow
\mathbb{R}^{1000}.
\]

Nel nostro progetto abbiamo invece soltanto due classi:
\textit{fake} e \textit{real}.

Per questo motivo sostituiamo l'ultimo layer:

\begin{verbatim}
num_features = self.model.fc.in_features
self.model.fc = nn.Linear(num_features, 2)
\end{verbatim}

L'attributo \texttt{in\_features} restituisce il numero
di caratteristiche in ingresso al layer originale,
che nel caso di ResNet18 è pari a 512.

La funzione \texttt{nn.Linear(512, 2)} costruisce
un nuovo layer lineare che trasforma le 512 caratteristiche
in due output.

Matematicamente, indicando con
$h(x)\in\mathbb{R}^{512}$ il vettore delle
caratteristiche estratte dalla ResNet18, il nuovo
classificatore calcola:

\[
z(x)=Wh(x)+b,
\]

dove:

\[
W\in\mathbb{R}^{2\times512},
\qquad
b\in\mathbb{R}^{2},
\qquad
z(x)\in\mathbb{R}^{2}.
\]

I due elementi di $z(x)$ rappresentano gli score
associati alle classi \textit{fake} e \textit{real}.

Il numero di parametri del nuovo layer è:

\[
512\cdot2+2=1026,
\]

ottenuto sommando i 1024 pesi della matrice $W$
e i due termini di bias.

\begin{importantbox}{Quali pesi sono pre-addestrati?}
Quando sostituiamo il layer finale, manteniamo i
pesi pre-addestrati di tutti i layer precedenti,
mentre i 1026 parametri del nuovo classificatore
vengono inizializzati casualmente.

Di conseguenza, il modello dispone già di un sistema
di estrazione delle caratteristiche visive, ma deve
ancora imparare a utilizzare queste caratteristiche
per distinguere immagini reali e generate.
\end{importantbox}

\subsubsection{Forward pass}

Il metodo \texttt{forward()} definisce come il modello
elabora gli input:

\begin{verbatim}
def forward(self, x):
    return self.model(x)
\end{verbatim}

Quando richiamiamo:

\begin{verbatim}
outputs = model(images)
\end{verbatim}

PyTorch esegue automaticamente il metodo
\texttt{forward()}, applicando la ResNet18
alle immagini fornite in input.

Considerando il mini-batch $b$ con $m_b$ immagini,
l'input ha dimensioni:

\[
X^{(b)}
\in
\mathbb{R}^{m_b\times3\times224\times224}.
\]

Il modello restituisce:

\[
Z^{(b)}
=
f_{\theta}(X^{(b)})
\in
\mathbb{R}^{m_b\times2}.
\]

Ogni riga di $Z^{(b)}$ contiene i due score
relativi a una singola immagine.

\begin{definitionbox}{Logits e probabilità}
Gli output prodotti dal classificatore sono chiamati
\textbf{logits}. Non rappresentano direttamente
probabilità, poiché non sono vincolati all'intervallo
$[0,1]$ e la loro somma non deve necessariamente
essere pari a uno.

Per ottenere le probabilità delle due classi,
possiamo applicare la funzione \textit{softmax}:

\[
p_k(x)
=
\frac{\exp(z_k(x))}
{\displaystyle\sum_{j=0}^{1}\exp(z_j(x))},
\qquad k\in\{0,1\}.
\]

Nel nostro dataset:

\[
p_0(x)=P_{\theta}(\text{fake}\mid x),
\qquad
p_1(x)=P_{\theta}(\text{real}\mid x).
\]

La classe predetta corrisponde all'indice del
logit maggiore:

\[
\hat{y}(x)
=
\arg\max_{k\in\{0,1\}} z_k(x).
\]
\end{definitionbox}

\subsubsection{Parametri addestrabili}

Un ultimo aspetto importante riguarda la differenza
tra caricare un modello pre-addestrato e congelarne
i parametri.

Nel nostro codice, i parametri della ResNet18
non vengono congelati. Pertanto, per impostazione
predefinita, sia i layer pre-addestrati sia il
nuovo classificatore finale risultano addestrabili.

Durante il training potremo scegliere tra due approcci:

\textbf{Feature extraction}: mantenere fissi i
parametri dei layer pre-addestrati e aggiornare
soltanto quelli del classificatore finale.

\textbf{Fine-tuning}: aggiornare anche i parametri
pre-addestrati, adattando progressivamente
l'intera rete al nuovo task.

\begin{warningbox}{Stato attuale del modello}
Il file \texttt{models.py} definisce e inizializza
il classificatore, ma non esegue alcun addestramento
sul dataset real/fake.

Il nuovo layer finale ha ancora pesi casuali,
quindi le predizioni non sono necessariamente
significative per il nostro problema.

Il training sarà implementato separatamente
nel file \texttt{train.py}, dove definiremo
la funzione di perdita, l'ottimizzatore
e il ciclo di aggiornamento dei parametri.
\end{warningbox}



\subsection{\texttt{train.py}}

Il file \texttt{train.py} implementa l'addestramento del
classificatore ResNet18 sul dataset di immagini reali e generate.

Il primo esperimento utilizza una strategia di
\textbf{feature extraction}: i parametri dell'estrattore
di caratteristiche pre-addestrato su ImageNet vengono
mantenuti fissi, mentre vengono ottimizzati soltanto
i parametri del nuovo classificatore finale.

Il training viene effettuato mediante mini-batch,
utilizzando la Cross-Entropy come funzione di perdita
e Adam come algoritmo di ottimizzazione.

\subsubsection{Congelamento dei parametri}

Il modello viene inizializzato mediante:

\begin{verbatim}
model = ResNet18Classifier().to(DEVICE)
\end{verbatim}

La chiamata \texttt{to(DEVICE)} trasferisce i parametri
e i buffer del modello sul dispositivo specificato,
che nel nostro caso può essere CPU, CUDA oppure MPS.

Successivamente vengono congelati tutti i parametri:

\begin{verbatim}
for parameter in model.parameters():
    parameter.requires_grad = False
\end{verbatim}

L'attributo \texttt{requires\_grad} indica se PyTorch
deve registrare le operazioni necessarie a calcolare
le derivate rispetto a un determinato parametro.
Impostandolo a \texttt{False}, il parametro non
riceverà gradienti durante il backward pass.
Riattiviamo quindi il calcolo dei gradienti
soltanto per il classificatore finale:

\begin{verbatim}
for parameter in model.model.fc.parameters():
    parameter.requires_grad = True
\end{verbatim}

Il modello possiede complessivamente
$11\,177\,538$ parametri, dei quali soltanto
$1\,026$ risultano addestrabili.

\begin{importantbox}{Congelamento dei parametri}
Impostare \texttt{requires\_grad=False} non equivale
a rimuovere i parametri dal modello.

I pesi pre-addestrati continuano a essere utilizzati
durante il forward pass, ma non vengono aggiornati
dall'ottimizzatore.

Il nostro problema consiste quindi nell'apprendere
un classificatore lineare sulle 512 caratteristiche
estratte dalla ResNet18.
\end{importantbox}

\subsubsection{Una particolarità: eval() durante il training}

Nel nostro codice, dopo aver congelato i parametri,
troviamo:

\begin{verbatim}
model.eval()
\end{verbatim}

Normalmente, durante l'addestramento di una rete
neurale, si utilizza:

\begin{verbatim}
model.train()
\end{verbatim}

mentre \texttt{model.eval()} viene utilizzato
durante la validazione e il test.

Nel nostro caso scegliamo deliberatamente di
utilizzare \texttt{eval()} anche durante il training.

Il motivo riguarda principalmente i layer
di \textbf{Batch Normalization} presenti in ResNet18.

Durante il training standard, questi layer
normalizzano le attivazioni utilizzando statistiche
calcolate sul mini-batch e aggiornano delle
statistiche interne, dette \textit{running statistics}.

In modalità valutazione, invece, utilizzano
le statistiche accumulate durante l'addestramento
precedente, senza aggiornarle.

\begin{definitionbox}{Batch Normalization}
Una trasformazione di Batch Normalization può essere
espressa, per una singola caratteristica, come:

\[
\operatorname{BN}(x)
=
\gamma
\frac{x-\mu}{\sqrt{\sigma^2+\varepsilon}}
+\beta,
\]

dove $\gamma$ e $\beta$ sono parametri
apprendibili, mentre $\mu$ e $\sigma^2$
sono le statistiche utilizzate per normalizzare
le attivazioni.

In modalità training, per la normalizzazione
vengono utilizzate le statistiche del mini-batch.

In modalità evaluation, vengono utilizzate
le stime memorizzate nelle running statistics.
\end{definitionbox}

Nel nostro esperimento vogliamo mantenere fisso
l'intero estrattore di caratteristiche, comprese
le sue statistiche di Batch Normalization.
Per questo motivo manteniamo il modello in modalità
\texttt{eval()} anche durante l'ottimizzazione.

\begin{warningbox}{eval() non disabilita i gradienti}
La chiamata \texttt{model.eval()} modifica il
comportamento di alcuni layer, come Batch Normalization
e Dropout, ma \textbf{non disabilita il calcolo dei gradienti}.

Di conseguenza, nonostante il modello sia in modalità
evaluation, i 1026 parametri del classificatore
continuano a ricevere gradienti durante
\texttt{loss.backward()}.

La disabilitazione del calcolo automatico dei gradienti
è invece controllata da \texttt{torch.no\_grad()}.

Questa scelta è appropriata per il nostro primo
esperimento di feature extraction, ma non rappresenta
la configurazione standard di un fine-tuning completo.
\end{warningbox}

\subsubsection{Funzione di perdita e ottimizzatore}

La funzione di perdita viene definita tramite:

\begin{verbatim}
criterion = nn.CrossEntropyLoss()
\end{verbatim}

Per ogni osservazione $i$, la rete restituisce
due logit:

\[
z_i=(z_{i,0},z_{i,1}),
\]

corrispondenti alle classi \textit{fake} e \textit{real}.

La Cross-Entropy per il mini-batch $b$ è:

\[
\mathcal{L}_b(\phi)
=
-\frac{1}{m_b}
\sum_{i\in b}
\log
\left(
\frac{\exp(z_{i,y_i})}
{\sum_{k=0}^{1}\exp(z_{i,k})}
\right).
\]

L'ottimizzatore viene inizializzato come:

\begin{verbatim}
optimizer = torch.optim.Adam(
    model.model.fc.parameters(),
    lr=LEARNING_RATE
)
\end{verbatim}

L'argomento passato all'ottimizzatore è importante:

\begin{verbatim}
model.model.fc.parameters()
\end{verbatim}

In questo modo Adam riceve esclusivamente
i parametri del classificatore finale e non
quelli dell'intera rete.

Il learning rate è specificato tramite
\texttt{LEARNING\_RATE}, definito in
\texttt{config.py}.

\subsubsection{Backward pass e aggiornamento dei parametri}

Per ogni mini-batch viene eseguita la seguente
sequenza di operazioni:

\begin{verbatim}
optimizer.zero_grad()

outputs = model(images)
loss = criterion(outputs, labels)

loss.backward()
optimizer.step()
\end{verbatim}

Le due istruzioni più importanti da approfondire
sono \texttt{zero\_grad()} e \texttt{backward()}.

\paragraph{Differenziazione automatica}

La chiamata:

\begin{verbatim}
loss.backward()
\end{verbatim}

esegue la differenziazione automatica della
funzione di perdita rispetto ai parametri
per i quali è attivo \texttt{requires\_grad}.

Nel nostro caso vengono calcolate le componenti:

\[
\nabla_\phi\mathcal{L}_b(\phi)
=
\left(
\frac{\partial\mathcal{L}_b}{\partial\phi_1},
\ldots,
\frac{\partial\mathcal{L}_b}{\partial\phi_{1026}}
\right).
\]

La successiva chiamata:

\begin{verbatim}
optimizer.step()
\end{verbatim}

utilizza tali gradienti per aggiornare i parametri
del classificatore secondo l'algoritmo Adam.

\subsubsection{Calcolo delle metriche di training}

Durante ogni epoca vengono accumulate la perdita
e il numero di classificazioni corrette.

Per la loss utilizziamo:

\begin{verbatim}
batch_size = labels.size(0)
train_loss += loss.item() * batch_size
\end{verbatim}

È importante osservare che, con la configurazione
predefinita di \texttt{CrossEntropyLoss},
\texttt{loss.item()} restituisce la perdita
\textbf{media del mini-batch}.

Per questo la moltiplichiamo per
\texttt{batch\_size}: vogliamo recuperare la
somma delle perdite delle osservazioni del batch.

Al termine dell'epoca calcoliamo:

\begin{verbatim}
train_loss /= train_total
\end{verbatim}

ottenendo:

\[
\overline{\mathcal{L}}_{\mathrm{train}}
=
\frac{
\displaystyle\sum_{b=1}^{B}
m_b\mathcal{L}_b
}{
\displaystyle\sum_{b=1}^{B}m_b
}.
\]

La ponderazione con $m_b$ è necessaria perché
l'ultimo mini-batch potrebbe contenere meno
di 32 osservazioni.

Se calcolassimo semplicemente la media delle
loss dei mini-batch, assegneremmo lo stesso peso
a batch di dimensioni differenti.

\paragraph{Calcolo dell'accuracy}

Le predizioni vengono ottenute tramite:

\begin{verbatim}
predictions = outputs.argmax(dim=1)
\end{verbatim}

La funzione \texttt{argmax(dim=1)} seleziona,
per ciascuna osservazione, l'indice del logit
più elevato.

Poiché i logit hanno dimensioni
$(m_b,2)$, l'operazione restituisce
un vettore di $m_b$ etichette predette.

Il numero di classificazioni corrette viene
calcolato con:

\begin{verbatim}
train_correct += (
    predictions == labels
).sum().item()
\end{verbatim}

L'accuracy finale è:

\[
\operatorname{Accuracy}_{\mathrm{train}}
=
\frac{1}{n_{\mathrm{train}}}
\sum_{i=1}^{n_{\mathrm{train}}}
\mathbb{I}(\hat y_i=y_i).
\]

\begin{warningbox}{Interpretazione della training loss}
La training loss viene accumulata durante l'epoca,
mentre i parametri del classificatore cambiano
dopo ogni mini-batch.

Pertanto, la loss di training riportata
non coincide necessariamente con la perdita
dell'intero training set valutata utilizzando
i parametri finali dell'epoca.

La validation loss, invece, viene calcolata
con il modello fissato al termine dell'epoca.
\end{warningbox}

\subsubsection{Validazione senza aggiornamento dei pesi}

Al termine di ogni epoca, il modello viene
valutato sul validation set.

Il codice utilizza:

\begin{verbatim}
with torch.no_grad():
    for images, labels in val_loader:
        outputs = model(images)
        loss = criterion(outputs, labels)
\end{verbatim}

Il contesto \texttt{torch.no\_grad()} disabilita
la registrazione delle operazioni per la
differenziazione automatica.

Questo riduce il consumo di memoria e il lavoro
necessario durante la validazione.

Inoltre, non vengono eseguite le chiamate:

\begin{verbatim}
loss.backward()
optimizer.step()
\end{verbatim}

Di conseguenza, durante la validazione non vengono
aggiornati i parametri del classificatore.

\begin{importantbox}{Differenza tra eval e nograd}
\texttt{model.eval()} imposta il comportamento
dei layer che distinguono tra training ed evaluation.

\texttt{torch.no\_grad()} disabilita la costruzione
del grafo computazionale per il calcolo dei gradienti.

Le due istruzioni non sono equivalenti.

Nel nostro esperimento utilizziamo
\texttt{eval()} anche durante il training,
ma \texttt{no\_grad()} soltanto durante
la validazione.
\end{importantbox}

\begin{importantbox}{Best checkpoint}
Il modello salvato è quello che ha ottenuto
la minore validation loss tra le epoche eseguite:

\[
t^*
=
\arg\min_{t\in\{1,\ldots,T\}}
\mathcal{L}_{\mathrm{val}}(\phi_t).
\]

Il checkpoint contiene lo stato del modello
all'epoca $t^*$ e non necessariamente quello
dell'ultima epoca.

Questa selezione permette di evitare di scegliere
automaticamente il modello dopo l'ultimo
aggiornamento, che potrebbe aver peggiorato
le prestazioni sui dati di validazione.

Per riutilizzare il checkpoint sarà necessario
ricostruire la stessa architettura e caricare
i pesi mediante \texttt{load\_state\_dict()}.

Il salvataggio del solo \texttt{state\_dict()}
non include lo stato dell'ottimizzatore:
è sufficiente per eseguire successive predizioni,
ma non per riprendere esattamente il training
dal punto in cui era stato interrotto.
\end{importantbox}

\begin{warningbox}{Utilizzo del test set}
La selezione del modello attraverso la validation loss
costituisce già una forma di scelta basata sui dati.

Per stimare le prestazioni del classificatore
selezionato su osservazioni mai utilizzate
per l'addestramento o per la selezione,
dovremo valutarlo successivamente sul test set.

Il test set non deve quindi influenzare
la scelta del numero di epoche o del checkpoint.
\end{warningbox}


\subsection{\texttt{evaluate.py}}

Valutazione del classificatore tramite metriche e predizioni sul test set.

\subsection{\texttt{\_\_init\_\_.py}}

File che identificano le cartelle del codice come package Python.


\section{Spiegazioni controfattuali: \texttt{src/explainers/}}

\subsection{\texttt{cf.py}}

Esperimenti con \textit{Counterfactual Instances}.

\subsection{\texttt{cem.py}}

Esperimenti con \textit{Contrastive Explanation Method}.

\subsection{\texttt{cfproto.py}}

Esperimenti con \textit{Counterfactuals Guided by Prototypes}.

\subsection{\texttt{cfrl.py}}

Esperimenti con \textit{Counterfactuals with Reinforcement Learning}.


\section{Esperimenti, dati e documentazione}

\subsection{\texttt{experiments/01\_exploration.ipynb}}

Notebook per verificare ambiente Python, configurazione, caricamento dati
e primi test dei modelli.

\subsection{\texttt{data/real/} e \texttt{data/fake/}}

Immagini delle due classi. FFHQ è previsto per i volti reali; occorre
aggiungere una sorgente per quelli generati.

\subsection{\texttt{checkpoints/}}

Pesi dei modelli salvati durante e dopo l'addestramento.

\subsection{\texttt{results/}}

Metriche, grafici ed esempi di immagini controfattuali.

\subsection{\texttt{slides/presentation.tex}}

Presentazione per il ricevimento intermedio e l'esposizione finale.

\subsection{\texttt{docs/project\_reference.tex}}

Questo documento di consultazione, da aggiornare durante il progetto.


\section{Stato del lavoro}

\begin{examplebox}{Configurazione iniziale}
\begin{itemize}
    \item Repository GitHub creata e struttura iniziale predisposta.
    \item Ambiente \texttt{XAIvenv} configurato sul Mac; PyTorch e MPS verificati.
    \item Accesso al server UniFi verificato, da utilizzare per esperimenti
          computazionalmente più impegnativi.
    \item \texttt{config.py} e \texttt{dataset.py}: prime versioni impostate;
          esecuzione e contenuto effettivo da verificare in VS Code.
    \item Notebook di esplorazione previsto con celle di controllo
          dell'ambiente e del dataset.
\end{itemize}
\end{examplebox}

\begin{importantbox}{Prossimo traguardo}
Preparare un piccolo dataset \textit{real/fake}, verificare i DataLoader
e caricare ResNet18. Solo dopo iniziare il confronto dei metodi Alibi.
\end{importantbox}

% =================================================
\chapter{Appunti tecnici}
% =================================================

% Aggiungere qui spiegazioni e risultati, man mano che vengono verificati.

\section{Dataset e preprocessing}
\textit{Da completare: sorgenti dati, classi, split, trasformazioni e controlli.}

\section{Classificatore ResNet18}
\textit{Da completare: architettura, fine-tuning, iperparametri e metriche.}

\section{Metodi controfattuali di Alibi}
\textit{Da completare: prerequisiti, compatibilità e risultati dei quattro metodi.}

\section{Esperimenti e risultati}
\textit{Da completare: impostazione sperimentale, tabelle e osservazioni.}

\section{Decisioni e questioni aperte}
\textit{Da aggiornare dopo i test e il ricevimento con il professore.}

\end{document}